\question\TFQuestion{F}{A population will always grow smaller if immigration is greater than death rate.}
\question\TFQuestion{F}{If negative feedback occurs when \textit{N} is greater than \textit{K}, then positive feedback occurs when \textit{N} is less than \textit{K}.}
\question\TFQuestion{F}{Negative feedback most often occurs in a population when \textit{N} is much less than \textit{K}.}
\question\TFQuestion{T}{Negative feedback most often occurs in a population when \textit{N} is much larger than \textit{K}.}
\question\TFQuestion{T}{Disease killed off most of the sea urchins around Jamaican coral reefs.  The number of species on the coral reef decreased dramatically after the sea urchins died. The sea urchins may have been a keystone species.}
\question\TFQuestion{F}{Any one trophic level supports on average about 10\% of the biomass of the next higher trophic level.}
\question\TFQuestion{F}{A field containing 30,000 kg of plant biomass will support about 300 kg of tertiary consumer biomass.}
\question\TFQuestion{F}{The carrying capacity of source habitats decreases when the carrying capacity of sink patches increases.}
\question\TFQuestion{T}{The carrying capacity of a habitat may change as environmental conditions change.}
\question\TFQuestion{F}{Evenness measures the distribution of relative abundance across all species in a population.}
\question\TFQuestion{T}{The brightly colored pattern of bees and wasps is an example of aposomatic coloration. }
\question\TFQuestion{F}{Heterochrony can be defined as the retention of juvenile characteristics in the adult.}
\question\TFQuestion{T}{Mutualism can be thought of as two species that ecologically exploit each other over the long term.}
\question\TFQuestion{F}{Competitive exclusion will occur when two species use even one identical resource out of many resources.}
\question\TFQuestion{F}{If populations of a single species are geographically isolated for \textit{any} period of time, then they will evolve biological reproductive isolation.}